\section{Open Challenges and Future Directions}
\label{sec:challenges}

Despite rapid progress, significant challenges remain. We organize these into technical, software engineering, and human factors categories.

\subsection{Technical Challenges}

\subsubsection{Reliability and Consistency}

Current agents remain unreliable on complex tasks. Specific issues include:
\begin{itemize}
    \item \textbf{Long-horizon failure}: Success rates decrease dramatically with task complexity
    \item \textbf{Non-determinism}: Same inputs can produce different outputs
    \item \textbf{Error cascades}: Early mistakes compound through multi-step tasks
    \item \textbf{Debugging difficulty}: Agent reasoning is often opaque
\end{itemize}

\subsubsection{Context and Scalability}

Large codebases present fundamental challenges:
\begin{itemize}
    \item Context windows, while growing, still cannot hold entire large projects
    \item Retrieval quality varies significantly across codebases
    \item Cross-repository dependencies are poorly handled
    \item Real-time context updates during editing remain difficult
\end{itemize}

\subsubsection{Verification and Testing}

Agents struggle to verify their own work:
\begin{itemize}
    \item Test generation quality lags code generation
    \item Agents cannot reliably assess whether changes are correct
    \item Formal verification integration remains limited
\end{itemize}

\subsection{Software Engineering Challenges}

\subsubsection{Code Quality}

Generated code may be functional but problematic:
\begin{itemize}
    \item Maintainability and readability concerns
    \item Inconsistent adherence to project style guides
    \item Violation of design patterns and idioms
    \item Technical debt accumulation from AI-generated code
\end{itemize}

\subsubsection{Security}

Security risks are significant:
\begin{itemize}
    \item Vulnerability introduction (common weakness patterns in training data)
    \item Secret leakage in prompts and context
    \item Supply chain risks from suggested dependencies
    \item Insufficient sandboxing of agent execution
\end{itemize}

\subsubsection{Intellectual Property}

Legal and ethical concerns persist:
\begin{itemize}
    \item License compliance in generated code
    \item Attribution requirements
    \item Copyright status of AI-generated code
    \item Training data provenance questions
\end{itemize}

\subsection{Human Factors}

\subsubsection{Skill Development}

Impact on developer learning is unclear:
\begin{itemize}
    \item Junior developers may miss foundational learning
    \item Over-reliance risks atrophy of debugging skills
    \item New skills required (prompt engineering, agent supervision)
    \item Educational curriculum implications
\end{itemize}

\subsubsection{Workflow Integration}

Practical integration challenges:
\begin{itemize}
    \item Code review of AI-generated code requires new approaches
    \item Team collaboration patterns with AI tools
    \item Trust calibration---when to accept vs. verify suggestions
    \item Organizational policies and governance
\end{itemize}

\subsection{Research Directions}

We identify several promising research directions:

\textbf{Better benchmarks}: Evaluation that captures real-world software engineering beyond bug fixes, including feature development, refactoring, and documentation.

\textbf{Human-agent collaboration}: Understanding optimal interaction patterns, trust calibration, and division of labor between humans and agents.

\textbf{Explainability}: Making agent reasoning transparent and debuggable.

\textbf{Domain-specific agents}: Specialized agents for security, testing, DevOps, and specific frameworks.

\textbf{Multi-agent systems}: Coordinating multiple specialized agents for complex software development tasks.

\textbf{Longitudinal studies}: Understanding long-term impact on code quality, developer skills, and team dynamics.
